\begin{proposition}[The exact surface-pairing obstruction in dimension three]
\label{c:threefold-pairing-kernel}
\status{Proved reduction; the remaining primitive pairing is unresolved}
Let $\mathcal X$ be a projective strictly semistable model of a smooth
projective threefold $X$ over a mixed-characteristic local field with
finite residue field $k$, and let $\ell\ne\operatorname{char}k$.
Write $Y^{(r)}$ for the disjoint union, over $\bar k$, of the
$r$-fold intersections of distinct components of its special fibre.
All cohomology below has coefficients $\Q_\ell$.  Put
\[
 A=H^2(Y^{(1)}),\quad U=H^2(Y^{(2)}),\quad
 B=H^0(Y^{(3)})(-1),\quad Z=H^0(Y^{(4)})(-1).
\]
Use the signed restriction and Gysin maps of the weight spectral
sequence, and define
\[
 \begin{split}
 R&=\operatorname{im}\bigl(\rho_A:A\longrightarrow U\bigr),\\
 K_0&=R\cap\gamma_B
       \ker\bigl(\rho_B:B\longrightarrow Z\bigr),\\
 b(u,v)&=\sum_{S\subset Y^{(2)}}\operatorname{Tr}_S(u_Sv_S)
           \ \in\Q_\ell(-2).
 \end{split}
\]
Then $K_0\subset\operatorname{rad}(b|_R)$, and there is a
Galois-equivariant isomorphism
\begin{equation}
\label{eq:c-threefold-radical-kernel}
 \ker\!\left(
 N:\operatorname{Gr}^W_4H^3(X_{\bar K})(1)
       \longrightarrow\operatorname{Gr}^W_2H^3(X_{\bar K})
 \right)
 \simeq
 \frac{\operatorname{rad}(b|_R)}{K_0}.
\end{equation}
Consequently the first remaining C.4 endpoint assertion is equivalent
to nondegeneracy of the induced form on $R/K_0$.

Choose a relatively ample class $L$.  The same assertion is
equivalent to nondegeneracy of the surface intersection form on
\begin{equation}
\label{eq:c-threefold-primitive-image}
 R_{\rm prim}
 =R\cap
   \ker\bigl(
       L:H^2(Y^{(2)})\longrightarrow H^4(Y^{(2)})(1)
       \bigr).
\end{equation}
In particular the remaining condition involves only a primitive
restriction image in the degree-two cohomology of the
double-intersection surfaces.  It does not require the unknown
nondegeneracy of an additional Gysin image.
\end{proposition}

\begin{proof}
Let
\[
 C=H^4(Y^{(1)})(1),\qquad D=H^2(Y^{(3)}).
\]
The weight spectral sequence of
\citet[Satz~2.10, formulas~(2.10.1)--(2.11)]{c_rapoportzink1982},
with the same normalization as Theorem~\ref{c:threefold-weight-gate},
gives
\[
 \begin{split}
 \operatorname{Gr}^W_4H^3(1)
  &=
  \frac{\ker\{U\oplus Z\longrightarrow C\oplus D,\ 
      (u,z)\longmapsto
      (\gamma_Uu,\rho_Uu+\gamma_Zz)\}}
       {\{(\gamma_Bb,\rho_Bb):b\in B\}},\\
 \operatorname{Gr}^W_2H^3
  &=
  \frac{\ker\{U\oplus Z\longrightarrow D,\
      (u,z)\longmapsto\rho_Uu+\gamma_Zz\}}
       {\{(\rho_Aa+\gamma_Bb,\rho_Bb):
                    a\in A,\ b\in B\}}.
 \end{split}
\]
Here $\gamma_U:U\to C$ and $\gamma_Z:Z\to D$ are the
signed codimension-one Gysin maps.
These displayed identifications use degeneration at $E_2$ by
the different Frobenius weights.  The monodromy map is induced
by the identity of $U\oplus Z$.  We choose the common
restriction/Gysin sign convention in both rows, as is possible
because the differential commutes with $N$.

If a class lies in its kernel, its representative has the form
$(\rho_Aa+\gamma_Bb,\rho_Bb)$.  Subtracting the boundary
$(\gamma_Bb,\rho_Bb)$ leaves $(r,0)$ with $r=\rho_Aa\in R$.
The remaining cocycle conditions are $\gamma_Ur=0$ and
$\rho_Ur=0$.  The latter is automatic from $\rho_U\rho_A=0$.
Two such representatives differ by a domain boundary precisely
when their difference is $\gamma_Bb$ with $\rho_Bb=0$.
Thus the kernel is
\[
 (R\cap\ker\gamma_U)/
       (R\cap\gamma_B\ker\rho_B).
 \]
Poincare duality identifies $\gamma_U$ with the adjoint of
$\rho_A$ under the perfect pairing between $A$ and $C$ and
the pairing $b$ on $U$.  Hence
$\ker\gamma_U=R^\perp$.
The relation $\gamma_U\gamma_B=0$ also shows that the denominator
lies in the radical.  This proves
\eqref{eq:c-threefold-radical-kernel}, including its Galois
equivariance and all its twists.
The equivalence with the C.4 endpoint follows from
Theorem~\ref{c:threefold-weight-gate}.

For the primitive formulation we apply
\citet[Theorem~0.2, complete proofs in \S\S1.1--2.5,
and the dimension-three converse in \S2.7,
author pp.~2--9]{c_msaito2003positivity}.
His dimension-three criterion consists of
\eqref{eq:c-threefold-primitive-image} and the primitive Gysin
image in $H^3(Y^{(1)})$.
The latter condition is automatic: \S2.7 identifies it with
the degree-two monodromy isomorphism, already proved by
restriction to a general hyperplane surface and
weight--monodromy for surfaces.
The other degree-zero and degree-one pairing conditions are
settled by rational positivity and Picard-variety
polarizations, as in the proof of
Theorem~\ref{c:improved-zero-weight}.
The remaining canonical pairing is exactly the intersection
form on the primitive degree-two part of each surface,
up to a common nonzero sign.  This proves the last equivalence.
\end{proof}

\begin{corollary}[Removing the rational divisor restriction image]
\label{c:threefold-divisor-split}
\status{Proved unconditional splitting; the spectral identification uses surface Tate}
Retain the notation of Proposition~\ref{c:threefold-pairing-kernel}.
Let $N_r=\operatorname{NS}(Y^{(r)})_\Q$ for $r=1,2$,
and let $\rho_{\rm NS}:N_1\to N_2$ be signed restriction.
Define the rational space
\[
 D_{\rm prim}=
 \operatorname{im}(\rho_{\rm NS})\cap
 \ker\!\left(
 N_2\longrightarrow
 \bigoplus_{S\subset Y^{(2)}}\Q,\quad
 (d_S)_S\longmapsto(\deg(L_Sd_S))_S
 \right).
\]
We regard $D_{\rm prim}\otimes\Q_\ell(-1)$ as a subspace of
$R_{\rm prim}$ by the divisor cycle map.  Then
\[
 R_{\rm prim}
 =\bigl(D_{\rm prim}\otimes\Q_\ell(-1)\bigr)
   \mathbin{\perp}R_{\rm rem},
 \qquad
 R_{\rm rem}=
 R_{\rm prim}\cap
       \bigl(D_{\rm prim}\otimes\Q_\ell(-1)\bigr)^\perp .
\]
The first summand has a nondegenerate negative-definite
rational intersection form.  Consequently the remaining
threefold assertion is equivalent to nondegeneracy of the
form on $R_{\rm rem}$.

If the divisor Tate conjecture holds for every smooth
projective surface over finite fields of the residue
characteristic, let
$T_r\subset H^2(Y^{(r)})$ be the Frobenius spectral summand
whose eigenvalues are not $|k|$ times roots of unity.
Then
\[
 R_{\rm rem}=\rho_A(T_1)\subset T_2.
\]
Thus, under that separate surface-Tate hypothesis, the
remaining pairing is entirely on the image of the
transcendental degree-two spectral summand.
\end{corollary}

\begin{proof}
All the divisor classes and restriction maps descend to a
finite extension of $k$.  The displayed primitive condition
is a rational matrix of divisor intersection numbers.
Kernels and images commute with extension from $\Q$ to
$\Q_\ell$, so its realized space is exactly the primitive
part of the image of rational divisor classes.
The surface Hodge index theorem makes the intersection
form negative definite on this rational primitive space,
component by component; the same rationality step is
spelled out in
\citet[Claim~4.6 and its full proof, author pp.~12--13]
{c_ito2004threefolds}.
In particular its Gram determinant is nonzero in $\Q$ and
remains nonzero in $\Q_\ell$.
Orthogonal projection onto this subspace, using the inverse
Gram matrix, preserves $R_{\rm prim}$ because the subspace
lies in $R_{\rm prim}$.  Its kernel is $R_{\rm rem}$, which
proves the asserted orthogonal direct sum.  This argument
does not identify the subspace with every algebraic vector
of $R_{\rm prim}$.

For the conditional assertion, the surface-Tate hypothesis
implies the divisor Tate conjecture on all the smooth
projective strata and all their finite constant-field
extensions by
\citet[Theorem~8, pp.~89--90]{BKahn2023Divisors};
the same reduction is used in
Theorem~\ref{thm:B-relative-surface-gate}.
The divisor subspace of $H^2$ has a Galois-stable orthogonal
complement: use the ample Lefschetz intersection form,
whose restriction to divisors is nondegenerate.
The divisor subspace has eigenvalues $|k|$ times roots of
unity.  Its complement has none: otherwise after a finite
constant extension it would have a nonzero eigenvector of
eigenvalue $|k'|$, contradicting Tate surjectivity and the
direct sum.  This also excludes a generalized eigenspace
of that type, since every nonzero generalized eigenspace
contains an eigenvector.  No additional full Frobenius
semisimplicity is used.

Restriction therefore preserves the two spectral blocks.
On the surface target the second block is orthogonal to
every divisor class and hence is primitive, because the
ample class is a divisor.  We obtain
\[
 R_{\rm prim}
   =\bigl(D_{\rm prim}\otimes\Q_\ell(-1)\bigr)
        \mathbin{\perp}\rho_A(T_1).
\]
Comparing with the unconditional orthogonal decomposition
proves the formula for $R_{\rm rem}$.
\end{proof}

\begin{remark}[The geometric content not removed by the pairing criterion]
\label{c:threefold-pairing-scope}
\status{Proved scope clarification; the primitive nondegeneracy is open}
The radical formula is an exact identification of the existing
kernel.  The primitive criterion additionally removes the
nonprimitive Lefschetz pieces and the Gysin-image condition.
It does not prove the resulting nondegeneracy.
In particular, containment of a $\Q_\ell$-subspace in the span
of divisor classes does not allow one to apply rational
Hodge-index positivity to that subspace without establishing
its rational origin.  For example, over $\Q_5$, choose
$\iota^2=-1$ with $\iota\equiv2\pmod5$.
The line $\Q_5(1,\iota)$ is isotropic for $-x^2-y^2$ although
this form is negative definite over $\Q$.
It is stable under the scalar Frobenius $2I$.
This is a counterexample to that linear-algebra inference,
not a geometric counterexample to weight--monodromy.

For the same reason, an isomorphism of a stratum motive with
an abelian motive would not alone identify its geometric
primitive form with a positive polarization.  The conditional
image-projector and positivity inputs in
\citet[Theorem~0.4 and \S\S4.1--4.3]{c_msaito2003positivity}
remain separate hypotheses.  We have not silently supplied
either input from Frobenius purity or from the 2026 theorem
for CM abelian varieties.
The corresponding warning is explicit in
\citet[submitted version, Remark~4.10 and footnote~3,
p.~14]{DMilne2002Polarizations}: equality of motive categories
under a Tate hypothesis does not by itself transfer the
positivity of the geometric forms.  The published-version
Theorem~3.3 and Corollary~3.4 concern abelian varieties.
\end{remark}
